\documentclass[10pt]{article}

\usepackage[T1]{fontenc}
\usepackage[margin=1in]{geometry}
\usepackage{amsmath}
\usepackage{booktabs}
\usepackage{enumitem}
\usepackage{microtype}
\usepackage{xurl}
\usepackage[hidelinks]{hyperref}

\newcommand{\project}{\textsc{zkTH06}}
\newcommand{\thgame}{Touhou Koumakyou v1.02h}
\newcommand{\code}[1]{\texttt{#1}}

\title{\project: Toward Replay-Verified Zero-Knowledge Execution\\
of a Legacy Bullet-Hell Game}
\author{N0zoM1z0}
\date{Living draft --- 10 August 2026}

\begin{document}
\maketitle

\begin{abstract}
We investigate whether a complete run of a legacy bullet-hell game can be
expressed as a deterministic transition system and eventually proven inside a
zero-knowledge virtual machine.  The central difficulty is semantic, not merely
cryptographic: the proof must represent the behavior of the shipped 32-bit
Windows executable rather than a convenient but divergent port.  We develop a
strict replay parser, a deterministic Linux compatibility runner, and a
frame-level differential methodology for \thgame.  An experimental canonical
digest projection is emitted by the runtime and independently checked by C++
and Python.  Four locally pinned six-stage replays complete under the
compatibility runner, and repeated complete canonical runs produce
byte-identical traces.  An address-bound Wine experiment additionally matches
the shipped executable and Linux reference for 2,000 frames on a 47-field raw
projection, while observing ambient x87 control word \code{0x007f} at every
post-calc anchor.  A first \code{no\_std} player-position microkernel then uses
integer dyadics and explicit PC24 round-to-nearest-even rather than host
floating point; it matches 1,999 consecutive retail-derived transitions.  An
enclosing transition now derives life state, invulnerability timer, time-stop
carry, movement rate/bounds, bomb multipliers, and character speeds from the
preceding state and one fixed configuration instead of accepting eleven
per-frame environment words.  The same crate runs in an OpenVM v2.0.1 RV32IM
guest whose public value commits the private input masks and complete endpoint.
Its tracked application proof binds the exact guest executable.  A further
Player refinement derives focus, previous input, fire-timer cadence, and 1,590
\code{SpawnBullets} requests while keeping the 4,018-byte private payload to
replay masks.  Its 1,999-transition OpenVM proof executes 5,468,598
instructions at an estimated 211,737,944 cells, proves in 75.57 seconds with
approximately 46.8 GiB peak resident memory, and verifies in 0.11 seconds.  A
fourth guest refines 1,590 local Reimu-A callback invocations through the
80-slot allocator and initialized geometry for ranks 1--3.  It hashes all 422
initialized bullets, executes 7,394,181 instructions at 309,751,316 cells, and
produces a tracked proof in 98.60 seconds.  The independently observed pre-call
slot arrays are not yet linked across frames, so this remains a local function
proof rather than an enclosing bullet-subsystem proof.  A
hash-bound static audit attributes 5,980 x87
instructions to mapped game functions and identifies direct transcendental,
conversion-helper, store-boundary, and control-profile obligations.  These
are followed by a pinned x87/SoftFloat counterexample search in which
34,002,802 tested result/condition/exception tuples match.  A separate
temporary i386 harness executes the verified original \code{\_\_ftol2} body;
1,000,014 in-domain EDX:EAX results match exact dyadic truncation, and a
further 1,000,154 focused executions match the 25-way ECL-variable
classification.  These two audited families are materialized as a hash-bound,
byte-free ledger of 321 original addresses.  A derived source/sink ledger
identifies 68 omission candidates and nine retained calls, while keeping every
correspondence, slicing, and refinement status open.
For the eight retained score conversions, a further address-bound audit and
total rational/exceptional Lean model avoid an item-finiteness assumption:
infinities are separated by a finite player box, while NaN maps through the
modeled helper invalid path to low EAX zero.  A dependent 95-instruction
player-position audit and total Lean clamp theorem further show that either
infinity is clamped into the required center range and isolate NaN candidate
exclusion, writer completeness, scheduling, and code binding as the remaining
player obligations.  A fixed-data audit then walks 8,384 subroutine
instructions from the seven retail ECL files: none of 1,844 candidate outputs
names player position,
including 80 increment/decrement operations that bypass the readonly guard.
Extraction and runtime-decoder binding remain open.  A sealed Linux probe
observes 2,081 collected items without a range violation, but does not promote
that corpus result into a proof.
These results establish a first narrow retail anchor, fail-closed enclosing
transition profile, and arithmetic proof surface only; the fixed post-calc
anchor, collision/death/respawn, bombs, time-stop writers, exact x87
transcendental semantics, presentation-dependent state, and a complete
retail canonical-subsystem exporter remain open equivalence gates.  This living paper
records the research questions, implementation evidence, negative results, and
changes in proof boundary as the project develops.
\end{abstract}

\section{Introduction}

A zero-knowledge proof can establish that a computation was executed correctly
without revealing all of its inputs.  For a game, this suggests a statement of
the form
\[
  S_{t+1} = F(S_t, I_t),
\]
where $S_t$ is gameplay state, $I_t$ is a private frame input, and $F$ is the
game's transition function.  A proof could commit to the game version, initial
state, route, input stream, and terminal claim while keeping the frame inputs
private.

The equation hides the hard part: what exactly is $F$?  \thgame{} is a closed,
legacy executable whose reconstructed source mixes gameplay with animation,
effects, message scripts, floating-point library behavior, and global random
number consumption.  A deterministic rewrite is not automatically equivalent
to the executable.  Proving the rewrite would only make its differences
cryptographically undeniable.

\project{} therefore places differential semantics before proof engineering.
The immediate objective is a replay runner that can expose the first differing
field at the first differing frame between an instrumented original or exact
reference and a portable implementation.  Only after representative full
replays agree do we freeze a proof-oriented kernel and select a zkVM backend.

Our intended claim is deliberately narrow.  It proves that an input sequence is
a valid execution under a committed game transition and reaches an explicit
terminal predicate.  It does not prove that a human generated the inputs, that
no external tool was used, or that a replay file is authentic.

\section{Background and Source Boundary}

The primary executable specification is the matching-oriented reconstruction
maintained by GensokyoClub~\cite{gensokyoclub-th06}.  It targets the Japanese
v1.02h executable and is valuable precisely because it is organized around
comparison with that binary.  It is not an official source release.  The
portable Linux/headless adaptation used by this project is derived from a
separate snapshot~\cite{n0zom1z0-headless}.  \project{} imports only a curated
source snapshot under \path{reference/}; upstream history, CI, reverse-engineering
automation, and unrelated documentation are not the project history.

The target executable is bound by SHA-256:
\begin{center}
\small\ttfamily
9f76483c46256804792399296619c1274363c31cd8f1775fafb55106fb852245
\end{center}
The commercial executable and DAT archives are not redistributed.  A
machine-readable manifest records every required filename and digest, and a
local verifier checks a user-supplied installation before experiments.

\subsection{Replay semantics}

A v1.02h replay contains character and difficulty metadata, per-stage state
snapshots, and run-length encoded input changes.  The shipped playback path
restores lives, bombs, power, rank, score-related counters, and RNG seed at stage
boundaries.  Its reversible byte transform and additive checksum detect some
accidental corruption but provide no authentication.

Consequently, ``the replay player reaches the ending'' is not a sound proof
statement.  A malicious witness can forge favorable stage snapshots.  A
canonical proof must derive stage $i+1$ from the live end state of stage $i$ and
either ignore the supplied snapshot or constrain it to equality.

\subsection{The ECL virtual machine is not the game}

Enemy ECL is a genuine domain-specific virtual machine with program counters,
integer and floating-point registers, control flow, call stacks, movement, and
bullet opcodes.  However, ECL reads player position, difficulty, rank, and enemy
state; it consumes the global RNG and mutates bullets, items, GUI state, and
score.  An isolated ECL proof would therefore accept an unconstrained
environment.  ECL is one subsystem of the frame transition, not the complete
proof boundary.

\section{Research Questions}

\begin{description}[style=nextline,leftmargin=0pt]
  \item[RQ1: Semantic boundary.]
  What canonical state projection is sufficient for observational equivalence
  with v1.02h gameplay across all reachable states in the evaluation corpus?

  \item[RQ2: Differential diagnosis.]
  How can heterogeneous executions---the original 32-bit Windows program and a
  portable Linux reconstruction---emit a stable, versioned representation that
  identifies the first divergent subsystem and field?

  \item[RQ3: Untrusted replay witnesses.]
  How can the compact input information in \code{.rpy} files be reused while
  preventing their per-stage checkpoints from becoming trusted state resets?

  \item[RQ4: Semantics-preserving slicing.]
  Which rendering, audio, animation, message, and effect operations can be
  removed, and which must remain as dependency-preserving shadow semantics due
  to timing, control flow, collision geometry, or shared RNG consumption?

  \item[RQ5: Arithmetic and slicing soundness.]
  Can the shipped program's x87, binary32 storage, and transcendental behavior
  be represented exactly at practical proof cost; and what machine-checked
  refinement justifies removing state or substituting integer arithmetic?

  \item[RQ6: Proof decomposition.]
  For each explicitly bounded slice, and after full semantic equivalence and
  refinement soundness for broader claims, which state layout, arithmetic
  model, frame chunking, and recursive aggregation strategy makes a spell,
  stage, and full six-stage run practical in a zkVM?
\end{description}

RQ1--RQ5 remain prerequisites for a whole-game answer to RQ6.  The OpenVM
experiments select a backend only for a retail-anchored Player slice.
The first proof hash-binds per-frame environment records; the second derives
those records inside a fail-closed enclosing player-state profile; and the
third derives focus, previous input, fire-timer cadence, and the
\code{SpawnBullets} request schedule.  All commit only fixed configuration,
replay masks, and their reached endpoints after the first experiment.  None
preselects the final whole-game architecture.  Corpus agreement may motivate a
projection or arithmetic design, but cannot by itself discharge RQ5's
universal claim.

\section{Methodology}

\subsection{Fail-closed replay ingestion}

The replay loader owns the decoded buffer and rejects invalid magic, version,
checksum, character/shot metadata, difficulty, unaligned or non-monotonic stage
offsets, malformed block sizes, unsupported input bits, regressing frame
numbers, and missing or malformed playback sentinels.  Playback cursors are
bounded by the first validated sentinel rather than by ambient process memory.

\subsection{Authentic frame order}

Replay input is injected by the existing replay callback at calculation
priority 5, after the game manager and before Stage, Player, Enemy/ECL,
Effects, Bullets/Items, and GUI.  Dialogue can restart the calculation chain;
this behavior is retained.  Complete Shoot, Bomb, Focus, Skip, and direction
masks are replayed.  Compatibility mode follows deaths, bombs, and respawns
instead of applying the reinforcement-learning harness's first-hit terminal.

\subsection{Two-level canonical tracing}

Writing every field of hundreds of entity slots on every frame would generate
multi-gigabyte traces.  The differential design therefore has two levels:

\begin{enumerate}
  \item a fixed-size per-frame record containing metadata and cryptographic
  digests for ordered subsystems; and
  \item a field-level canonical snapshot for a selected tick around the first
  digest mismatch.
\end{enumerate}

Every serializer uses explicit little-endian integer widths, raw IEEE-754 bit
patterns, stable array slot indices, and offsets relative to loaded script
bases.  It never hashes pointers, compiler padding, locale-dependent text, or
unordered container iteration.  Subsystem separation distinguishes global/RNG,
Player, player bullets, Enemy/ECL, enemy bullets, lasers, items, Stage,
GUI/message, ANM shadow, and Effects.

\subsection{Differential loop}

For each pinned replay, both runners must use the same data hashes and replay
bytes.  The comparison proceeds in frame order, first comparing the root digest,
then subsystem digests, then a rerun with the field-level snapshot at the first
failing tick.  A final score match is supporting evidence but never substitutes
for frame equivalence.

\section{Implementation Status}

The public repository uses a flattened, project-owned history.  The imported
engine is isolated under \path{reference/}; its pinned revisions and license
boundary are explicit.  Changes made for \project{} remain ordinary reviewable
commits on top of that snapshot.

As of 9 August 2026, the following components are implemented:

\begin{itemize}
  \item a standalone Python replay inspector and a stricter owning C++ loader;
  \item direct headless compatibility playback from the first populated stage;
  \item bounded, complete-mask input injection through the original callback;
  \item deterministic JSON diagnostics and explicit terminal reasons;
  \item a retail-data manifest and local SHA-256 verifier;
  \item Linux CI that builds the reference runner and cross-checks both replay
  validators on three valid and three deliberately malformed cases.
\end{itemize}

The current compatibility path intentionally restores each replay stage
snapshot because that mirrors shipped playback and enables oracle development.
It is not yet canonical proof playback.  No zkVM guest is frozen or claimed to
be implemented.

\section{Preliminary Evaluation}

Four six-stage v1.02h replays are pinned by hash.  Three are tracked fixtures
from a public replay uploader; one world-record replay remains local because its
source prohibits secondary redistribution.  All four reached the
\code{replay-complete} terminal in the Linux compatibility runner.

\begin{table}[h]
  \centering
  \small
  \begin{tabular}{@{}llrrr@{}}
    \toprule
    Replay & Route & Logic ticks & Stage 6 frame & Final score \\
    \midrule
    \code{th6\_ud000134} & Lunatic / Marisa B & 84,817 & 16,485 & 181,778,020 \\
    \code{th6\_ud000232} & Lunatic / Reimu B  & 86,676 & 16,727 & 136,978,560 \\
    \code{th6\_ud002677} & Normal / Reimu A, NN/NB & 85,759 & 17,283 & 172,519,700 \\
    \code{th6\_udLuRB}   & Lunatic / Reimu B  & 118,605 & 26,233 & 804,515,970 \\
    \bottomrule
  \end{tabular}
  \caption{Compatibility playback results.  NN/NB denotes no-miss,
  no-bomb.  These are Linux self-consistency measurements, not original-game
  equivalence results.}
  \label{tab:compatibility-results}
\end{table}

Two independent 1,000-tick runs of \code{th6\_ud002677} produced
byte-identical JSONL with SHA-256
\begin{center}
\small\ttfamily
d4606beb09ea9040b0e747d95b2dc7ebe70d36f5b5c152c8289b6074e69e899a.
\end{center}
On the initial host, full untraced runs completed in approximately
1.1--1.6 seconds and stayed below 40 MiB resident memory.  These timings are
environment-specific and are not a performance claim.

The Stage 6 final score is computed during gameplay: playback restores the
preceding stage's score at Stage 6 entry but does not restore Stage 6's own end
score.  Exact score agreement is therefore a useful end-to-end consistency
check.  It remains much weaker than frame-level state equivalence.

\subsection{Canonical projection experiment}

Two independent complete revision-0.2 traces of the Normal Reimu A NN/NB run
were byte-identical across all 85,759 frame records.  Both terminated at Stage
6 frame 17,283 with score 172,519,700.  Each trace occupied 50,769,392 bytes and
had SHA-256
\begin{center}
\small\ttfamily
ea0cdf948ba7668cba31064dfa421a9c279fdab28bdbd1d57c817ba13db84117.
\end{center}
Although only fixed-size digests were retained, each run hashed
6,192,210,130 bytes (5.767 GiB) of selected canonical payload.  Two concurrent
local executions took 31.90 and 32.04 seconds and peaked at 38,388 and 38,724
KiB resident memory.  This measures the current diagnostic implementation; it
does not estimate zkVM proving cost.

An earlier serializer hashed every named member of each Stage ANM VM.  Two
2,000-tick runs then disagreed in the Stage subsystem at record zero, while RNG
and all other subsystem digests agreed.  Inspection showed that Stage VMs are
allocated with \code{malloc} and that disabled interpolation members are not
initialized.  Encoding those members only when their control mode can read
them restored byte identity for both the 2,000-tick and complete runs.  This is
a concrete example of why a canonical state must follow future-read semantics
rather than C++ object layout.

The complementary inactive-state failure is machine-checked in a small Lean
model.  Two free Effect slots with equal active-only projections but distinct
raw radius bits evolve to distinct projected states under the same spell-orbit
allocation.  A second theorem shows that retaining dormant radius and angle
makes this modeled allocation phase commute.  The proof compiles with Lean
4.32.0 without \code{sorry}, admitted results, or custom axioms.  Its scope is
only this allocation phase and does not establish sufficiency for all Effect
callbacks or bind the model to compiled C++.

\subsection{Static arithmetic census}

The arithmetic audit binds the executable, authoritative function map, and
implemented-function list by SHA-256.  GNU \code{objdump} 2.38 linearly decoded
133,078 instruction lines in the complete image, including 10,100 x87 lines.
Of these, 9,522 lie in mapped ranges.  Complete-image counts are static upper
bounds: they include CRT, D3DX, and potentially unreachable code rather than
dynamic execution frequencies.

The mapping attributes 5,980 x87 instructions to 159 \code{th06::} functions,
all marked implemented by the pinned reconstruction metadata.  Table~\ref{tab:arithmetic-census}
summarizes the source-binding candidate.  The 26 direct transcendental sites
are all \code{fsincos}.  The store widths show why rounding every source
operation to binary32 would also be wrong: x87 intermediates can remain live
until a specific DWORD or QWORD store.

\begin{table}[h]
  \centering
  \small
  \begin{tabular}{@{}lr@{}}
    \toprule
    Static property in mapped game functions & Count \\
    \midrule
    x87 instructions & 5,980 \\
    direct x87 control/environment instructions & 0 \\
    \code{fsincos} instructions & 26 \\
    \code{frndint} instructions & 10 \\
    x87 stores & 1,793 \\
    DWORD stores & 1,713 \\
    QWORD stores & 80 \\
    \bottomrule
  \end{tabular}
  \caption{Static x87 census for mapped \code{th06::} ranges.  Counts are
  sites, not execution frequencies or a proved gameplay slice.}
  \label{tab:arithmetic-census}
\end{table}

Function-body counts alone are incomplete.  The audit also finds 163 direct
calls from mapped game code into 18 x87 helpers or wrappers: 77 call
\code{\_\_ftol2}, 16 each call \code{sin} and \code{cos}, 11 call
\code{sqrt}, and seven call \code{atan2}.  The complete image contains 356
XMM-operand instruction lines, none inside a mapped \code{th06::} range; they
belong to a Pentium-4 \code{pow} implementation, a CRT helper used by D3DX
rasterization, and feature detection.  Their removal still requires a
transitive noninterference argument for the corresponding draw paths.

The CRT trigonometric wrappers embed raw x87 control word \code{0x027f}.  A
small Lean definition checks its 53-bit significand precision,
round-to-nearest-even mode, and six masked exception bits.  Static inspection
did not find a game-body instruction that establishes this word at entry, so
the result is deliberately a profile fact rather than an environment theorem.
An original exporter must record and constrain the entry word and processor
profile before exact transcendental equivalence can be claimed.

\subsection{Arithmetic-boundary oracle experiment}

The arithmetic experiment compares x87 and Berkeley SoftFloat Release 3e for
add, subtract, multiply, divide, and square root under the decoded PC53
nearest-even profile.  It also compares binary32/binary64 \code{fstp},
\code{frndint}, and signed 32/64-bit \code{fistp} under nearest-even and the
audited D3DX toward-zero profile.  The runner exports exact upstream commit
\code{f74b1e48110ac3a27dd49b787d164e55e42d81d1}, verifies selected file hashes,
and compiles only the 32 required translation units with the SoftFloat
\code{8086} specialization.  The probe source for this run had SHA-256
\begin{center}
\small\ttfamily
cf55d2ed7bde6c1020877b427eca43bc0f4e49f8ab7ffa29bfc4d8fdb487b392.
\end{center}

Raw result representations and x87 exception bits 0--5 are compared.  The
SoftFloat flags map invalid, divide-by-zero, overflow, underflow, and inexact;
the x87-only denormal-operand bit is supplied by an instruction-specific rule.
The first such rule---set the bit whenever either arithmetic input is
subnormal---failed on a fixed subnormal-dividend/zero-divisor case.  Following
the documented exception priority and instruction-specific \code{frndint},
\code{fistp}, and \code{fstp} behavior repaired the model before the reported
campaign.  A finite-class Lean definition checks both counterexamples and the
boundary distinction, but does not prove that the manuals, hardware, or
SoftFloat implement that definition.

On an AMD EPYC 9654 with GCC 11.4.0, one million deterministic pseudorandom
cases per operation, profile, and input family, plus fixed boundary cases,
yielded Table~\ref{tab:softfloat-probe}.

\begin{table}[h]
  \centering
  \small
  \begin{tabular}{@{}lrrr@{}}
    \toprule
    Family & binary32-derived & PC53 extended & Mismatches \\
    \midrule
    Five basic operations & 5,000,798 & 5,001,314 & 0 \\
    Five boundaries, two RC profiles & 10,000,140 & 10,000,330 & 0 \\
    \bottomrule
  \end{tabular}
  \caption{Host-specific x87/SoftFloat result/exception counterexample search.}
  \label{tab:softfloat-probe}
\end{table}

All six exception bits occurred in the basic campaign: 1,001,228 invalid,
88,466 denormal-operand, 56 divide-by-zero, 255,746 overflow, 265,276
underflow, and 7,547,221 inexact observations.  Counts overlap when an
instruction sets multiple bits.  Across both campaigns, 30,002,582 tuples
matched and the measured run took 4.3 seconds.

Canonical PC53 inputs retain the full extended exponent range, set the
explicit integer bit consistently with the exponent, and clear the eleven
significand bits below the selected precision.  The zero-mismatch result makes
SoftFloat a plausible oracle for this subset on the measured host; it proves
neither implementation correct.  Condition codes, TOP/tag state, arbitrary
NaN inputs, \code{fld m32fp} denormal signaling, comparisons, remainder,
complete helper ABIs, transcendental instructions, dynamic reachability, and
model-to-code binding remain outside the experiment.

\section{Evolution of the Research Direction}

This section is intentionally chronological.  Superseded ideas remain visible
so that later design choices can be audited rather than reconstructed from
memory.

\paragraph{Phase A: prove the ECL virtual machine.}
The initial idea treated enemy ECL as the natural zkVM target.  ECL has an
explicit instruction format, registers, comparison state, call stack, program
counter, and interpreter loop, making it attractive for an opcode-constraint
prototype.

\paragraph{Phase B: recognize the open environment.}
Source inspection showed that ECL reads Player and GameManager state, shares the
global RNG, and mutates bullets, items, GUI, and scoring state.  A proof of ECL
alone could accept forged player positions or rank.  The claim was narrowed
from ``a TH06 pattern was executed'' to ``ECL executed under supplied external
inputs,'' which was judged insufficient as the main result.

\paragraph{Phase C: reframe the game as a frame transition.}
The project moved from inventing a larger VM to extracting a deterministic
gameplay kernel $S_{t+1}=F(S_t,I_t)$.  ECL became one ordered subsystem.  This
reframing also made replay inputs a natural witness serialization.

\paragraph{Phase D: distrust replay checkpoints.}
Inspection of ReplayManager revealed that each stage carries mutable resources
and RNG state that playback restores directly.  Merely proving successful
playback would allow forged resources at every boundary.  Canonical proof mode
must derive stage continuity and constrain or discard those snapshots.

\paragraph{Phase E: differential validation before ZK.}
The portable headless runner was found deterministic but not known equivalent:
it substitutes host math for x87 behavior and skips draw callbacks that can
write later-observable state.  The immediate milestone changed from circuit
construction to an exact replay differential runner against the shipped or an
exactly matched reference executable.

\paragraph{Phase F: digest first, snapshot on divergence.}
A naive full-state dump for 640 enemy bullets, 513 items, effects, enemies, and
ANM state on every frame is unnecessarily large.  The current design emits
ordered subsystem digests each frame, then reruns only the first mismatching
tick with a versioned field-level snapshot.  Stable slot identity and iteration
order are retained because allocation, collision, and RNG behavior can depend
on them.

Future changes to the central claim or state boundary should be appended here
with the evidence that motivated them.

\section{Limitations and Open Threats to Validity}

\begin{itemize}
  \item \textbf{Partial retail equivalence only.}  A shipped-executable Wine
  probe now agrees with the Linux reference for 2,000 frames and a 34-field
  raw projection on one replay.  It does not yet export the eleven canonical
  subsystem payloads, cover later stages or the replay matrix, prove that its
  input/timing interventions are noninterfering, or establish whole-state and
  whole-program equivalence.
  \item \textbf{The first microkernel is narrow.}  The integer PC24
  player-position transition matches 1,999 consecutive retail-derived steps,
  but it excludes dead/spawning writes and exceptional or subnormal arithmetic,
  and it does not model hitbox/orb/animation side effects.  Its dyadic rounding
  algorithm, reachable-domain checks, interior x87 control state, retail
  instruction binding, Rust/zkVM code binding, and projection noninterference
  are not yet machine-checked.
  \item \textbf{Floating point.}  The target uses x87 instructions including
  trigonometric and rounding operations, explicit DWORD/QWORD stores, and
  conversion helpers.  The static audit identifies wrapper word
  \code{0x027f}, while the retail probe observes ambient \code{0x007f} at all
  2,000 post-calc boundaries.  Neither result proves entry state or every
  transient per-instruction word, and the current PC53 SoftFloat campaign does
  not cover the observed PC24 ambient profile.  Hardware
  transcendental approximations also require a processor/emulator profile.
  Host \code{libm}, spill behavior, and intermediate precision may diverge by
  one or more bits and later alter collisions.  A pinned SoftFloat experiment
  agrees with hardware results and six exception bits for tested arithmetic,
  store, rounding, integer-conversion, and comparison cases, but it is neither
  formally verified nor complete x87 state semantics.  The comparison probe
  includes C0--C3 and selected NaNs, while a separate exact-byte
  \code{\_\_ftol2} probe covers full register results only for canonical finite
  signed-i64 inputs.  Load exceptions, arbitrary/noncanonical NaNs, remainder,
  exceptional helper inputs, site-specific reachable ranges/observations,
  address binding, and transcendental behavior remain open.  A focused ECL
  campaign tests exceptional/default classification but is not universal
  semantics.  A hash-bound ledger indexes 321 comparison/helper addresses; a
  derived manual
  source ledger proposes 68 omissions and nine retained calls, but every
  correspondence, slice disposition, and helper refinement is still open.  Its
  bounded scan also leaves EDX live at 42 stopping points.  The ECL quotient
  still lacks verified decoding, universal helper/classifier agreement,
  jump-target binding, and guest refinement; the eight point-score calls lack
  discharge of the dependent player-position artifact, collision-code binding,
  architectural masked-invalid/helper-path refinement, and caller binding,
  together with reachable difficulty and pre-score bounds.  The 95-instruction
  player audit and total clamp theorem reduce its invariant to non-NaN movement
  candidates plus initialization, writer-completeness, scheduling, binary/x87,
  correspondence, and guest obligations; they do not prove those premises.  A
  fixed-data ECL audit finds zero player-position outputs among 1,844 candidates
  and checks all 80 operations using the two handlers that bypass the readonly
  guard, but still lacks a
  proof of archive extraction, runtime decoding/dispatch, instruction
  immutability, non-ECL alias completeness, and artifact-to-Lean binding.
  The total score model covers finite, infinite, and NaN item coordinates
  conditionally; 2,081 Linux samples exercise only its finite branch and prove
  none of those bindings.  The complete address-indexed semantics, code
  binding, and any fixed-point refinement have not yet been implemented or
  proved.
  \item \textbf{Draw-chain effects.}  Headless mode skips drawing, but some draw
  callbacks write flags, ANM fields, laser positions, or other state that may be
  read later.  Each removal needs dependency analysis and corpus validation.
  \item \textbf{Incomplete corpus.}  The current samples cover Normal and
  Lunatic full runs but not all characters, shot types, Extra, retries, and
  adversarial stage-checkpoint mutations.
  \item \textbf{Reconstruction gaps.}  The matching project is highly complete
  but not identical in every function.  Unmatched code must be closed or
  explicitly excluded from the accepted statement.
  \item \textbf{Claim scope.}  A valid execution proof is not a proof of human
  play, authorship, or the absence of tool assistance.
  \item \textbf{Projection soundness.}  Equal selected-field digests imply
  equality only of that projection under the SHA-256 collision-resistance
  assumption.  Until coverage is frozen and justified, omitted state can still
  invalidate an observational-equivalence claim.  In particular, inactive
  Effect slots retain future-live values, and dynamic screen-effect calc jobs
  are not yet represented even though screen shake consumes the shared RNG.
  \item \textbf{Initializedness is semantic.}  Some reconstructed structures
  contain disabled members left uninitialized by the reference allocator.
  Hashing or constraining such bytes creates false distinctions and may import
  undefined host state into the statement.  Current ANM guards are tested but
  not yet justified by an exhaustive future-read proof.
  \item \textbf{Model-to-code binding.}  The planned Lean refinement is not yet
  connected to the authoritative C++ transition or a zkVM guest.  Differential
  agreement can expose mistakes in such a connection but cannot prove it.
\end{itemize}

The commercial game data is verified locally by digest and is not
redistributed.  Replay fixtures retain source links and redistribution status;
inputs whose sources prohibit reuse remain local.

\section{Conclusion}

\project{} treats semantic equivalence as the first-class obstacle to proving
a legacy game's execution.  The initial replay runner demonstrates that the
reconstructed engine can consume real v1.02h witnesses deterministically and
complete full runs, while strict parsing removes several ambient-memory and
format-trust hazards.  The result is not yet a proof system and is not yet an
exact model of the shipped executable.

An experimental shared binary protocol now emits real per-frame subsystem and
root digests with raw floating-point bits, stable identities, and fail-closed
decoding.  Complete repeated Linux runs are byte-identical, and the first broad
ANM experiment already exposed the distinction between allocated memory and
future-live semantic state.  The next concrete milestone is to resolve the
known reuse and dynamic-job omissions, add selected field-level snapshots, and
emit the same schema from the shipped or an exact-reference executable.  This
should turn ``the replay desynchronized'' into a reproducible first-frame,
first-subsystem, first-field diagnosis.

That evidence will guide the projection, but it cannot prove a slice sound.
The intended end state additionally requires a machine-checked commuting
transition and noninterference argument, exact or refined arithmetic semantics,
and a reviewed binding from those definitions to both executables.  Only then
can zkVM performance results support the intended claim.


\bibliographystyle{plain}
\bibliography{references}

\end{document}
